\documentclass[journal]{IEEEtran}

\usepackage[utf8]{inputenc}
\usepackage{amsmath,amsfonts,amssymb}
\usepackage{booktabs}
\usepackage{graphicx}
\usepackage{cite}
\usepackage{url}
\usepackage{microtype}
\usepackage{multirow}
\usepackage{algorithm}
\usepackage{xcolor}
\usepackage{algpseudocode}
\graphicspath{{.}{figures/}{../figures/}}

\newcommand{\pending}[1]{\textcolor{red}{[PENDING: #1]}}
\newcommand{\code}[1]{\texttt{#1}}

\begin{document}

\title{Targeted Local Search, Not Adaptive Parameters: A Controlled Study of a Hybrid Genetic Algorithm for Curriculum-Based Course Timetabling}

\author{Amri~Reza~Wahyudin
        and~Rohmat~Indra~Borman%
\thanks{A. R. Wahyudin and R. I. Borman are with the Faculty of Engineering and Computer Science, Universitas Teknokrat Indonesia, Bandar Lampung 35142, Indonesia (e-mail: amri\_reza\_wahyudin@teknokrat.ac.id).}%
\thanks{Corresponding author: Amri Reza Wahyudin.}%
}

\markboth{Journal Submission, Vol.~XX, No.~X, 2026}%
{Wahyudin and Borman: Targeted Local Search, Not Adaptive Parameters}

\maketitle

\begin{abstract}
Hybrid genetic algorithms for curriculum-based course timetabling (CB-CTT) usually combine several mechanisms---adaptive parameter control, local search, repair operators---without isolating which of them produces feasible schedules. This paper presents a controlled study of a hybrid GA in which diversity-driven adaptive parameter control (APC) and violation-directed elite local search (ELS) are separated in a $2\times2$ factorial design augmented with a random-targeting control. On the 21 ITC~2007 CB-CTT instances, at a fixed budget of 50{,}000 objective evaluations, we find that ELS is the component that drives feasibility: arms without it never reach a hard-feasible schedule, whereas the directed-ELS arm solves every seed on \code{comp01} and \code{comp07}. Targeting is causal rather than incidental: with an identical budget and move set, violation-directed targeting reaches 10/10 feasible runs where random targeting reaches 0/10 on \code{comp07}. APC produces no measurable improvement over a static parameterisation at the tested base rate (GA\_ELS and HAGA differ by less than 3\% in soft-constraint penalty where both are feasible); because the adaptive arm's realised mean crossover rate (0.706) differs from the static 0.80, a mean-matched control is reported alongside the ablation. We report the complete comparison against simulated annealing, a canonical GA, and a deliberately naive discrete PSO, together with explicit gaps to published best-known solutions, and we release the protocol, manifests and analysis code. Across the full grid, simulated annealing attains the best median SCV among feasible runs (684.5) and the best mean rank on the search objective (1.700), while HAGA and GA\_ELS obtain feasible schedules far more often (77.3\% and 79.2\% versus 59.0\%); the Nemenyi test does not separate the three at $\alpha=0.05$, and the median gap to best-known solutions is 17.7$\times$, 22.3$\times$ and 22.5$\times$ respectively.
\end{abstract}

\begin{IEEEkeywords}
Curriculum-based course timetabling, genetic algorithms, adaptive parameter control, elite local search, controlled ablation, reproducibility.
\end{IEEEkeywords}

\section{Introduction}
\IEEEPARstart{U}{niversity} course timetabling is a constrained assignment problem in which lectures must be placed into rooms and timeslots while satisfying teacher, curriculum and room constraints. The curriculum-based variant (CB-CTT) was standardized by the Second International Timetabling Competition (ITC~2007)~\cite{mccollum2009setting,bonutti2012benchmarking}, and remains an active benchmark for metaheuristics~\cite{bettinelli2015survey,chen2021survey}.

Metaheuristics for CB-CTT rarely consist of a single operator. A typical hybrid combines a population method with local search, adds a repair or penalty mechanism, and tunes several control parameters. Reported gains are then attributed to the whole architecture, while the individual contributions remain confounded: the ablation changes the budget allocation, the base parameters, or both~\cite{eiben1999parameter}. The result is a literature in which it is difficult to tell \emph{which} mechanism is responsible for feasibility, and in which complex controllers are often added without a controlled test of their value.

This paper takes one such architecture and isolates its components. We study a hybrid GA with two mechanisms: (i) an adaptive parameter controller (APC) that maps population diversity to crossover and mutation rates, and (ii) a violation-directed elite local search (ELS). We ask a narrow question: \emph{which of the two is necessary, and does the direction of the local search matter?} The answer is obtained from a $2\times2$ factorial ablation with equal base parameters, plus a control arm in which ELS keeps its budget and move set but selects lectures uniformly at random.

Our contributions are:
\begin{enumerate}
    \item a controlled ablation of APC and ELS in which the base crossover rate is identical across cells, removing the confound present in earlier comparisons;
    \item a random-targeting control that isolates the \emph{direction} of ELS from its budget, showing that targeting is the causal ingredient for feasibility on the dense instance (10/10 versus 0/10) and improves SCV by 19\% on the small instance;
    \item evidence that APC has no measurable effect once base parameters are equalised, contrary to the usual assumption that diversity feedback is essential;
    \item a reproducible protocol with verified instance hashes, exact evaluation accounting, per-instance paired inference, and an explicit comparison against published best-known solutions.
\end{enumerate}

We do not claim state-of-the-art performance. On the contrary, the study is positioned as a mechanism isolation: the strongest baseline (simulated annealing) remains better on solution quality where both methods are feasible, and the gap to best-known solutions is reported explicitly. What the proposed hybrid offers is a different trade-off---more reliable feasibility on dense instances and comparable runtime (510 s versus 494 s) at an equal evaluation budget.

\section{Related Work}
\subsection{CB-CTT formulations and benchmarks}
The ITC~2007 formulation assigns each lecture to a room and timeslot subject to room occupancy, teacher and curriculum conflicts, and course unavailability; the objective penalizes room-capacity excess, violations of minimum working days, room instability and isolated curriculum lectures~\cite{mccollum2009setting}. The benchmark suite and its data format were documented by Bonutti et al.~\cite{bonutti2012benchmarking}, and best-known results are periodically consolidated in surveys~\cite{ceschia2022educational}. We follow the reference formulation and verify the 21 instance files against a published mirror (Section~\ref{sec:benchmark}).

\subsection{Metaheuristics for timetabling}
Simulated annealing, tabu search and hybrid evolutionary methods remain the dominant approaches~\cite{sylejmani2022simulated,rezaeipanah2021hybrid,lacerda2021systematic}. Competition results show that the best-performing solvers combine constructive initialisation with a local search that is aware of the constraint structure~\cite{song2021competition}. Population methods are usually paired with a local improvement operator; the relative contribution of the population and of the local search is rarely isolated.

\subsection{Adaptive parameter control}
Adaptive parameter control has a long history in evolutionary computation~\cite{eiben1999parameter,goldberg1989genetic}. Diversity-driven schedules are attractive because they are deterministic and cheap, but their benefit depends on whether the control law can actually traverse a meaningful range. A multiplicative law with a small slope, for example, remains within a few percent of its base value and is indistinguishable from a constant schedule. We make this point explicitly in Section~\ref{sec:method} and test it in Section~\ref{sec:results}.

\section{Problem Formulation}
\label{sec:problem}
Let $E$ be the set of lectures, $R$ the rooms and $T$ the timeslots. A schedule assigns each lecture $e\in E$ a pair $(t_e,r_e)\in T\times R$. Hard constraints are: at most one lecture per room and timeslot; no teacher teaches two lectures at the same time; no curriculum contains two lectures at the same time; and no lecture is placed in an unavailable period. Following the reference formulation, room-capacity excess is a \emph{soft} penalty (one point per student above capacity), together with minimum working days (5 points per missing day), room stability (1 point per additional room per course) and curriculum compactness (2 points per isolated lecture).

The search objective is the scalar
\begin{equation}
F(S) = W\cdot \mathrm{HCV}(S) + \mathrm{SCV}(S),
\end{equation}
where $\mathrm{HCV}$ counts hard violations, $\mathrm{SCV}$ is the official soft penalty, and $W=100{,}000$. A conservative bound is $\mathrm{SCV}\le 5\sum_c\max(0,m_c-1)+\sum_c(\ell_c-1)+2\sum_c \ell_c\,\mu_c$, where $m_c$ is the minimum working days of course $c$, $\ell_c$ its number of lectures and $\mu_c$ the number of curricula containing it. For all 21 instances this bound is below 25{,}600, so $F$ is lexicographic in $(\mathrm{HCV},\mathrm{SCV})$; a schedule is reported \emph{feasible} when $\mathrm{HCV}=0$.

\section{Proposed Method}
\label{sec:method}
\subsection{Encoding and diversity}
A schedule is a vector of $E$ integers, each encoding a room--timeslot pair, $c_e = t_e\cdot|R| + r_e$. Population diversity is the normalized mean pairwise Hamming distance
\begin{equation}
H_{\mathrm{norm}} = \frac{2}{N(N-1)}\sum_{i<j}\frac{1}{E}\sum_{k=1}^{E}\mathbb{1}\!\left[c_{i,k}\neq c_{j,k}\right],
\end{equation}
which is the average per-gene heterozygosity of the population. Its computation costs $O(N^2E)$ per generation; it is therefore computed only for arms that use it.

\subsection{Adaptive parameter control}
APC maps diversity to the crossover and mutation rates over their full admissible ranges:
\begin{align}
P_c &= P_{c,\min} + (P_{c,\max}-P_{c,\min})\,H_{\mathrm{norm}},\\
P_m &= P_{m,\min} + (P_{m,\max}-P_{m,\min})\,(1-H_{\mathrm{norm}}),
\end{align}
with $P_c\in[0.40,0.95]$ and $P_m\in[0.01,0.10]$. High diversity therefore favours crossover and suppresses mutation; low diversity does the opposite. The span form is deliberate: a multiplicative law $P_m=P_{m,\mathrm{base}}(1+\beta(1-H))$ with a small $\beta$ cannot leave a narrow band around its base value, which makes the controller indistinguishable from a constant. The arms \code{GA\_ELS} and \code{GA\_ELS\_Random} use static $P_c=0.80$, $P_m=0.05$; only \code{HAGA} and \code{GA\_APC} use the span law.

\subsection{Violation-directed elite local search}
At each generation, ELS is applied to the top $\lceil 0.10N\rceil$ individuals (Algorithm~\ref{alg:els}). The candidate set $\mathcal{K}$ is the set of lectures involved in hard violations; once the schedule is feasible, $\mathcal{K}$ is replaced by the set of lectures involved in soft violations (courses short of working days, unstable across rooms, or isolated inside a curriculum-day). A step selects $k\in\mathcal{K}$ and proposes either a reassignment to a uniformly random room--timeslot pair or a swap with another lecture (equal probability); the move is accepted only if $F$ improves. The pool is refreshed every ten steps. Swaps matter because a single-gene reassignment cannot exchange two lectures without passing through a duplicate-slot state.

\begin{algorithm}[!t]
\caption{Violation-Directed Elite Local Search}
\label{alg:els}
\begin{algorithmic}[1]
\For{each elite schedule $S$}
    \For{$\text{step}=1$ to $S_{\max}$}
        \If{$\text{step} \bmod 10 = 1$}
            \State $\mathcal{K}\gets$ lectures in hard violations of $S$
            \If{$\mathcal{K}=\emptyset$}
                \State $\mathcal{K}\gets$ lectures in soft violations of $S$
            \EndIf
        \EndIf
        \State select $k\in\mathcal{K}$ uniformly; $S'\gets$ reassign or swap $k$
        \State evaluate $F(S')$; deduct one evaluation
        \If{$F(S')<F(S)$} $S\gets S'$ \EndIf
    \EndFor
\EndFor
\end{algorithmic}
\end{algorithm}

\subsection{Two-stage budget allocation}
The evaluation budget is split explicitly: 70\% is spent on the generational population search and 30\% on a final intensification of the single best schedule, using the same move set and targeting rule. This prevents the common confound in which a memetic method spends most of its budget inside local search, and it gives the population phase and the intensification phase comparable, reportable shares.

\section{Experimental Design}
\label{sec:experiments}
\subsection{Instances and budget}
\label{sec:benchmark}
All 21 ITC~2007 CB-CTT instances are used; the files are byte-identical (SHA-256) to a published mirror of the ITC~2007 suite. Table~\ref{tab:instances} lists their sizes and best-known upper bounds~\cite{ceschia2022educational}. \code{comp15} is structurally identical to \code{comp03}; it is reported descriptively but excluded from all inferential tests, leaving $N=20$ unique instances. Every run uses exactly 50{,}000 objective evaluations, counted for offspring and local-search steps alike; 30 seeds (101--130) are used per instance and arm.

\begin{table}[!t]
\caption{Benchmark instances, sizes, and best-known upper bounds~\cite{ceschia2022educational}. $^\dagger$Structural duplicate of \code{comp03}; excluded from inference.}
\label{tab:instances}
\centering
\small
\begin{tabular}{l r r r}
\toprule
Instance & Lectures & Constraints & Best known \\
\midrule
\code{comp01} & 160 & 53 & 5 \\
\code{comp02} & 283 & 513 & 24 \\
\code{comp03} & 251 & 382 & 64 \\
\code{comp04} & 286 & 396 & 35 \\
\code{comp05} & 152 & 771 & 284 \\
\code{comp06} & 361 & 632 & 27 \\
\code{comp07} & 434 & 667 & 6 \\
\code{comp08} & 324 & 478 & 37 \\
\code{comp09} & 279 & 405 & 96 \\
\code{comp10} & 370 & 694 & 4 \\
\code{comp11} & 162 & 94 & 0 \\
\code{comp12} & 218 & 1{,}368 & 294 \\
\code{comp13} & 308 & 468 & 59 \\
\code{comp14} & 275 & 486 & 51 \\
\code{comp15}$^\dagger$ & 251 & 382 & 62 \\
\code{comp16} & 366 & 518 & 18 \\
\code{comp17} & 339 & 548 & 56 \\
\code{comp18} & 138 & 594 & 61 \\
\code{comp19} & 277 & 475 & 57 \\
\code{comp20} & 390 & 691 & 4 \\
\code{comp21} & 327 & 463 & 74 \\
\bottomrule
\end{tabular}
\end{table}

\subsection{Arms}
The main comparison uses five arms: (i) \code{Canonical\_GA} ($P_c=0.80$, $P_m=0.05$, no local search); (ii) \code{Discrete\_PSO}, a deliberately naive discrete recombination baseline in which each gene is copied from the personal best, the global best or a random slot with probabilities 0.40/0.40/0.20 and no velocity is used; (iii) \code{Simulated\_Annealing} ($T_0=100$, cooling 0.995, 50 steps per temperature, reassignment and swap moves); (iv) \code{GA\_ELS} (ELS, static parameters); and (v) \code{HAGA} (APC + ELS). The ablation adds \code{GA\_APC} (APC only) and \code{GA\_ELS\_Random} (ELS with random lecture selection, identical budget and moves), and is run on \code{comp01}, \code{comp07} and \code{comp12} with 10 seeds.

\subsection{Statistics}
The statistical unit is the instance; seeds are nested repetitions and are never pooled as independent samples~\cite{demsar2006statistical}. Medians across the 30 seeds form one observation per instance. We report the Friedman omnibus test with the Iman--Davenport correction, Kendall's $W$, and the Nemenyi critical difference at $\alpha=0.05$ and $0.10$. Per-instance paired Wilcoxon tests (HAGA versus each baseline) are Holm-corrected. Feasibility is reported with Wilson intervals.

\subsection{Reproducibility}
Each run stores its full chromosome, summary and monotone history; each experiment stores a manifest with instance hashes, source hashes, dependency versions and host metadata. Runs are resumable: a completed run is never overwritten, and re-invoking the same experiment identifier executes only the missing grid cells.

\section{Results}
\label{sec:results}
\subsection{Main comparison}
Table~\ref{tab:main} reports the five-arm comparison. We summarise the outcome in terms of three quantities: feasibility rate (HCV$=0$), median SCV among feasible runs, and median wall-clock runtime at equal evaluation budget. On the search objective (median fitness, which combines feasibility and soft penalties), simulated annealing ranks first (1.700), followed by GA\_ELS (2.075) and HAGA (2.225); the canonical GA and the naive PSO are last (4.000 and 5.000). The Friedman statistic is $\chi^2=65.33$ ($p=2.2\times10^{-13}$) with Iman--Davenport $F=84.6$ and Kendall's $W=0.817$. The Nemenyi critical difference is 1.364 at $\alpha=0.05$ and 1.230 at $\alpha=0.10$; the SA--HAGA rank gap (0.525) lies below both, so rank-based testing does not separate them. The arms differ in their outcome profile: HAGA and GA\_ELS reach feasibility in 77.3\% and 79.2\% of runs versus 59.0\% for SA, while SA's median SCV among feasible runs is 684.5 against 1{,}003.0 and 1{,}054.0. Per-instance paired tests on the objective, Holm-corrected, show HAGA significantly better than SA on \code{comp11} ($p=0.002$), significantly worse on \code{comp04}, \code{comp08}, \code{comp13} and \code{comp14} ($p<0.005$), and statistically indistinguishable on the remaining fifteen instances.

\begin{table*}[!t]
\caption{Main comparison across 21 ITC~2007 instances (30 seeds per arm, 50{,}000 evaluations). SCV is reported over feasible runs only; the gap is the median ratio to the best-known value of Table~\ref{tab:instances}. \code{comp11} has best-known $0$, so its ratio is undefined and excluded from the median.}
\label{tab:main}
\centering
\small
\begin{tabular}{l c c c c}
\toprule
Arm & Feasible runs & Median SCV (feasible) & Median runtime & Gap to best-known \\
\midrule
Canonical GA & 0/630 (0.0\%) & -- & 487.4 s & -- \\
Discrete PSO (naive)  & 0/630 (0.0\%) & -- & 497.9 s & -- \\
Simulated annealing   & 372/630 (59.0\%) & \textbf{684.5} & 494.1 s & \textbf{17.7}$\times$ \\
GA\_ELS               & \textbf{499/630 (79.2\%)} & 1{,}054.0 & 540.6 s & 22.5$\times$ \\
HAGA                  & 487/630 (77.3\%) & 1{,}003.0 & 509.8 s & 22.3$\times$ \\
\bottomrule
\end{tabular}
\end{table*}

\subsection{Controlled ablation}
Table~\ref{tab:ablation} isolates the components. Three findings are consistent across the ablation cells. First, ELS is necessary for feasibility: \code{Canonical\_GA} and \code{GA\_APC} never reach $\mathrm{HCV}=0$ on any seed, whereas \code{GA\_ELS} solves all ten seeds on both \code{comp01} and \code{comp07}, and \code{HAGA} solves 10/10 and 8/10 respectively. Second, targeting is the causal ingredient, not the presence of local search: with the same budget and move set, \code{GA\_ELS} solves 10/10 seeds on \code{comp07} while \code{GA\_ELS\_Random} solves 0/10. Third, APC adds no measurable benefit once the base parameters are equal: \code{GA\_ELS} and \code{HAGA} differ by 2.6\% on \code{comp01} (249.5 vs.\ 256.0) and by 2.0\% on \code{comp07} (2{,}462 vs.\ 2{,}413), in opposite directions. Because the adaptive law realises a mean $P_c$ of 0.706 (range 0.58--0.95), this cell compares adaptive control with a static rate of 0.80 rather than isolating adaptivity alone; a mean-matched static control ($P_c=0.70$, 10 seeds, \code{comp01}) is reported as median SCV 254.5 with 10/10 feasible runs, against 256.0 for the adaptive arm on the same seeds (a 0.6\% difference); at the original static rate of 0.80 the median is 249.5. The median effect of adaptivity is therefore indistinguishable at matched operating points, although the adaptive arm shows lower dispersion (IQR 66.8 versus 76.2 and 166.8), which $n=10$ cannot resolve.. Cells with a single feasible run (HAGA on \code{comp12}) are reported for completeness, not as estimates of central tendency.

\begin{table}[!t]
\caption{Controlled ablation (10 seeds per cell, 50{,}000 evaluations). Feasibility is HCV$=0$; SCV is the median over feasible runs.}
\label{tab:ablation}
\centering
\small
\begin{tabular}{ll c c}
\toprule
Instance & Arm & Feasible & Median SCV \\
\midrule
\multirow{5}{*}{\code{comp01}}
& Canonical GA      & 0/10 & -- \\
& GA\_APC           & 0/10 & -- \\
& GA\_ELS           & 10/10 & 249.5 \\
& GA\_ELS\_Random   & 10/10 & 307.0 \\
& HAGA              & 10/10 & 256.0 \\
\midrule
\multirow{5}{*}{\code{comp07}}
& Canonical GA      & 0/10 & -- \\
& GA\_APC           & 0/10 & -- \\
& GA\_ELS           & 10/10 & 2{,}462.0 \\
& GA\_ELS\_Random   & 0/10 & -- \\
& HAGA              & 8/10 & 2{,}413.0 \\
\midrule
\multirow{5}{*}{\code{comp12}}
& Canonical GA      & 0/10 & -- \\
& GA\_APC           & 0/10 & -- \\
& GA\_ELS           & 2/10 & 1{,}942.5 \\
& GA\_ELS\_Random   & 0/10 & -- \\
& HAGA              & 1/10 & 1{,}473.0 \\
\bottomrule
\end{tabular}
\end{table}

\subsection{Hyperparameter sensitivity}
Table~\ref{tab:sensitivity} reports the span-law sensitivity on \code{comp01}. All seven configurations reach 10/10 feasibility; the conditional SCV lies between 192.5 and 332.5. No configuration separates from the pre-registered baseline outside sampling noise, and the apparent advantage of a smaller population ($N=30$) is a trend on one instance, not a result we use for tuning. The table is reported as robustness evidence for the parameterisation used in the main comparison.

\begin{table}[!t]
\caption{Sensitivity of the span-law parameters on \code{comp01} (10 seeds, 50{,}000 evaluations). Feasibility is reported separately from the conditional SCV.}
\label{tab:sensitivity}
\centering
\small
\begin{tabular}{l c c c}
\toprule
Configuration & Feasible & Median SCV & SCV IQR \\
\midrule
Baseline ($N=50$, $P_{c,\min}=0.40$, $P_{m,\max}=0.10$) & 10/10 & 256.0 & 66.8 \\
$P_{c,\min}=0.30$ & 10/10 & 200.0 & 57.2 \\
$P_{c,\min}=0.55$ & 10/10 & 275.5 & 65.2 \\
$P_{m,\max}=0.05$ & 10/10 & 271.0 & 45.8 \\
$P_{m,\max}=0.20$ & 10/10 & 249.5 & 25.2 \\
$N=30$ & 10/10 & 192.5 & 50.0 \\
$N=100$ & 10/10 & 332.5 & 77.2 \\
\bottomrule
\end{tabular}
\end{table}

\subsection{Distance to best-known solutions}
Table~\ref{tab:instances} includes the best-known upper bound for each instance~\cite{ceschia2022educational}; those values were obtained with 300--500~s timeouts, which is comparable to the wall-clock of the runs reported here (about 500~s). The median gap to best-known (over instances with at least one feasible run, excluding \code{comp11} whose optimum is zero) is 17.7$\times$ for SA, 22.3$\times$ for HAGA and 22.5$\times$ for GA\_ELS; GA and PSO never produce a feasible schedule, so no gap is defined. \code{comp05} is infeasible for every arm under the 50{,}000-evaluation budget and \code{comp12} for most arms: for these two constraint-dense instances the budget is insufficient regardless of method. The gap is reported because it is the honest anchor for a method that does not claim state-of-the-art: it quantifies how far a fixed-budget, Python implementation remains from specialized solvers, and it prevents over-reading small differences among the arms.

\section{Discussion}
\label{sec:discussion}
\subsection{Why targeting matters}
The mechanism is visible in the neighbourhood. At a feasible schedule, the hard-conflict set is empty; an untargeted ELS therefore samples lectures uniformly and spends its budget on moves that cannot improve the objective. On a feasible \code{comp01} schedule, 0 of 400 random single-gene reassignments improved the composite fitness, while 9 of 400 swaps did. Directed targeting changes the candidate distribution exactly where the search would otherwise stall; random targeting leaves the budget intact but removes the direction, and the ablation shows the consequence (10/10 versus 0/10 on \code{comp07}).

\subsection{Why APC does not help}
The span law traverses $[0.40,0.95]$ and $[0.01,0.10]$, so the controller is no longer a near-constant, yet its effect remains within noise. Two reasons are plausible: selection pressure and ELS dominate the search, and the mutation range that matters is narrow once the schedule is feasible. The practical lesson is methodological: a controller must be tested against a static arm with the same base parameters, otherwise its effect cannot be separated from the parameter change that accompanies it.

\subsection{Budget allocation}
Splitting the budget 70/30 between population search and single-best intensification changed the result more than any parameter of the controller. This is consistent with the trajectory-based baselines, which spend the entire budget on one solution. We report the split explicitly so that future comparisons can reproduce it rather than infer it.

\section{Threats to Validity}
\textbf{Construct validity.} The scalar objective uses $W=100{,}000$; the bound in Section~\ref{sec:problem} shows that this is lexicographic for the 21 instances, but the constant is not derived per instance. Room capacity is soft, following the reference formulation; results are therefore not comparable to studies that treat it as hard.

\textbf{Internal validity.} Budget parity is defined in objective evaluations, not wall-clock time. The Python implementation and the $O(N^2E)$ diversity computation make wall-clock comparisons implementation-dependent; we report runtime but do not use it for inference. The SA baseline uses a geometric temperature schedule that the 50{,}000-evaluation budget truncates before the intended terminal temperature (0.67 rather than 0.001); combined with the $100{,}000$ hard-violation weight, this makes SA effectively greedy with respect to hard constraints---a conservative but strong baseline.

\textbf{External validity.} The study covers one problem family (CB-CTT) and one benchmark suite. \code{comp05} and \code{comp12}, the two densest instances, remain infeasible for every arm under the 50{,}000-evaluation budget; conclusions about feasibility should not be generalized to instances with substantially different constraint density.

\textbf{Conclusion validity.} Seeds are paired within an instance and never pooled across instances; per-instance tests use the seed as the pairing unit. \code{comp15} is a structural duplicate of \code{comp03} and is excluded from inference. Per-instance tests are Holm-corrected, and the Nemenyi critical difference is reported at two significance levels because the conclusions for adjacent ranks are sensitive to $\alpha$.

\section{Conclusion}
This paper separated the components of a hybrid genetic algorithm for CB-CTT. The controlled ablation shows that violation-directed elite local search is the component that produces feasible schedules, that its direction is causal (targeted 10/10 versus random 0/10 on \code{comp07}), and that diversity-driven parameter control adds no measurable effect once base parameters are equalised. The result is a narrower but more defensible claim than a superiority statement: the value of the architecture lies in targeted intensification under a fixed budget, not in the adaptivity of its parameters. Future work will test adaptive targeting rules, repair-based initialisation, and budget scaling against the best-known anchor.

\section*{Funding}
\pending{funding statement}

\section*{Conflict of Interest}
The authors declare no conflict of interest.

\section*{Data Availability}
The analysis scripts and manifests are included in the repository; the per-run records are available from the corresponding author on reasonable request. The aggregation and statistics reported here can be regenerated from the stored runs without re-executing the solvers.

\bibliographystyle{IEEEtran}
\bibliography{refs}

\end{document}
